\documentclass[10pt,a4paper]{amsart}
\usepackage[margin=19mm]{geometry}
\usepackage{amsmath,amssymb,mathtools}
\usepackage[scr=rsfs,cal=euler]{mathalfa}
\usepackage{xfrac}
\usepackage[unicode,hidelinks,bookmarks=false]{hyperref}
\newcommand{\ie}{i.e.\ }
\newcommand{\cf}{cf.\ }
\newcommand{\resp}{respectively\ }
\newcommand{\Subsection}{\subsection}
\providecommand{\nobreakdash}{\nobreak\hbox{-}}
\setlength{\emergencystretch}{2em}
\multlinegap0pt
\usepackage[shortlabels]{enumitem}
\usepackage{xcolor}
\usepackage{amssymb}
\newtheorem{lem}{Lemma}
\newcommand{\N}{\mathbb{N}}
\newcommand{\Z}{\mathbb{Z}}
\newcommand{\cM}{\mathcal{M}}
\newcommand{\cv}{\mathrm{conv}}
\title{Finitely Generated Congruences in Semirings and Canonical Positive Models}
\author{Snehinh Sen}
\date{Source: arXiv:2511.12815v2 (CC BY 4.0)}
\begin{document}
\maketitle

\noindent\emph{Source and licence.} Finitely Generated Congruences in Semirings and Canonical Positive Models. Version arXiv:2511.12815v2 (CC BY 4.0), distributed by arXiv under the Creative Commons Attribution 4.0 International licence, http://creativecommons.org/licenses/by/4.0/, as recorded in the arXiv licence metadata for this exact version and shown on the abstract page https://arxiv.org/abs/2511.12815v2. Full licence notice: \texttt{licenses/arxiv-2511.12815-CC-BY-4.0.txt}.\par\smallskip

\noindent\textit{Editorial scope.} Counted unit: the article's lem lem6 (source0.tex lines 814--817). The article's complete original proof of it is delivered with it (source0.tex lines 819--889, one \texttt{proof} environment of 71 lines). No mathematics is written, repaired or re-derived: the statement and the proof are the article's own lines, in the article's own notation, and no retained line is substituted or re-flowed.  Retained uncounted as the source-local context the proof needs: lines 774--782.  Nothing was omitted from any retained span, so the package carries no omission record.  No retained line contains a citation, so the wrapper carries no bibliography and no bibliography entry.  No retained line contains a figure environment, an image-inclusion command or drawing code, so no image is delivered and the package declares no asset dependency.  Editorial formatting only, in addition: an AMS \texttt{amsart} preamble that restates the article's own declarations and macros used by the retained lines, namely the environments \texttt{lem} on separate counters, together with the standard packages those lines need.  The retained environments print in this document as Lemma 1 in order of appearance, whereas the article prints the same items with the numbering of its own full document.  The complete native source of the article as distributed in the arXiv e-print for this exact version is bundled unchanged as \texttt{sources/189235/source0.tex}.
\begin{lem}
    \label{lem4}
    Given any $\gamma$-oriented collection $W=\{w_1,\ldots,w_n\}$ and a collection of indices $J\subseteq \{1,\ldots, n\}$ with $|J|\geq 2$ and any $\delta>0$, there is a $\gamma$-oriented collection $\{v_1,\ldots, v_n\}\in \cM(W)$ such that
    \begin{enumerate}
        \item $v_i=w_i$ if $i\notin J$,
        \item $|v_i|_\gamma < \delta $ for every $i\in J$, and
        \item $w_i \in \N\langle v_j | j\in J\rangle$ for each $i\in J$.
    \end{enumerate}
\end{lem}

\begin{lem}
    \label{lem6}
    Let $n\geq 3$. Given any $V = \{v_1,\ldots,v_n\} \in \cM(W)$ and any pair of indices $1\leq i,j\leq n$, such that $p(w_i)>p(w_j)$ there is a $U=\{u_1,\ldots,u_n\} \in \cM(V)$ such that $p(w_i-w_j) \in \cv(p(u_1),\ldots,p(u_n)).$
\end{lem}

\begin{proof}
    Without loss of generality, assume that $i=1, j=2$. As $V$ is a $\Z$-basis of $\Z^n$, there are integers $\mu_k$ such that $e =w_1-w_2 = \sum_k \mu_k v_k$. Let $\gamma_k=v_k\cdot \gamma >0$. So, 
    $$\sum_k \mu_k \gamma_k = e\cdot \gamma>0.$$
    Suppose $Q$ is the number of $1\leq k\leq n$ such that $\mu_k<0$.

\underline{\textbf{CASE I: $Q=0$}} \\
Just choose $U=V$.

\underline{\textbf{CASE II: $Q=1$}} \\
By possibly permuting the vectors, assume that $\mu_1 = -s<0$. Let $\delta = \frac{1}{sn}\min\{e\cdot \gamma, \gamma_1 \}>0$. By Lemma \ref{lem4}, we get a $U'=\{u_1',\ldots,u_n'\}\in \cM(V)$ such that $0<|u_j'|_\gamma<\delta$ for each $j\geq  2$ and $u_1'=v_1$. Furthermore, $v_k \in \mathrm{Span}_\N\langle u_j'|j\geq 2\rangle$ for $k\geq 2$. So, $w_1-w_2 =\sum_j \mu'_j
u_j'$ where $\mu_1'=\mu_1$ and $\mu_j'\geq 0$ if $j\geq 2$.

Consider the set of vectors $(\lambda_2,\ldots, \lambda_n)$ of non-negative integers such that 
\begin{enumerate}
    \item $\lambda_js \leq \mu'_j$ if $j\geq 2$
    \item If $\eta = \sum_{j\geq 2} \lambda_j u_j'$, then $\gamma_1>\eta\cdot \gamma$, and
\end{enumerate}

The set of such vectors $(\lambda_2, \ldots, \lambda_n)$ is finite and non-empty (as the zero vector satisfies these conditions). Choose a termwise maximal element of this set. Thus, for each $j\geq 2$, we must have either $\lambda_j = \left\lfloor\frac{\mu_j'}{s}\right\rfloor =:\lambda_j^{max}$ or $(\eta+u_j')\cdot \gamma > \gamma_1$. Otherwise, the vector $\lambda' = \lambda+e_{j}$, where $e_j$ is the $j$-th standard basis vector, satisfies all the three conditions, contradicting maximality.

Let $J_1 = \{2\leq j \leq n : \lambda_j=\lambda_j^{max}\}$ and $J_2 = \{2\leq j \leq n : j\notin J_1\}$.

Observe that 
\begin{align*}
\sum_{j\geq 2} \lambda_j^{max}u_j'\cdot \gamma
    &> \sum_{j\geq 2} \frac{\mu_j'u_j'}{s}\cdot \gamma - u_j'\cdot \gamma \\
    &= \frac{e\cdot \gamma}{s} + \gamma_1 - \sum_j u_j'\cdot \gamma \\
    &> \frac{e\cdot \gamma}{s} + \gamma_1 - n\delta > \gamma_1.
\end{align*}

So $J_2\neq \emptyset$. Also, for any $j\in J_2$, we have $\mu'_j \geq (\lambda_j+1)s$.

Define
$$u_j = \begin{cases}
    u_1' -\eta & j=1 \\
    u_j' & j\in J_1\\
    u_j'+\eta - u_1' & j\in J_2 \\
\end{cases}$$

It is immediate that $U =\{u_1,\ldots, u_n\} \in \cM(U')\subseteq \cM(V)$. Furthermore, 
$$w_1 - w_2 = \left(-s+\sum_{j\in J_2} (\mu'_j  - s\lambda_j)\right)  u_1 +\sum_{j\geq 2} (\mu'_j  - s\lambda_j)u_j$$
is in $\mathrm{Span}_\N(U)$. This proves our claim.

\underline{\textbf{CASE III: $Q\geq 2$}} \\
After possibly permuting the indices, suppose $\mu_1,\ldots, \mu_t<0$ and $\mu_{t+1},\ldots, \mu_n\geq 0$. Let $M=\sum_{k>t}\mu_k>0$ (as $e\cdot \gamma>0)$ and $\delta = {M}^{-1}(e\cdot \gamma)>0$. By Lemma \ref{lem4}, choose $U'=\{u_1',\ldots,u_n'\}\in \cM(V)$ such that $0<|u_j'|_\gamma<\delta$ for each $j\leq t$ and $u_j'=v_j$ if $j>t$. So, $e = \sum_k \mu_k'u'_k$, where $\mu_k'<0$ if $k\leq t$ and $\mu_k'=\mu_k$ for $k>t$.

We claim the existence of a $Z =\{z_1,\ldots,z_n\} \in \cM(U')$ such that if 
$e = \sum \theta_k z_k$, then $\theta_k=\mu_k'$ if $k\neq t$ and $\theta_t > 0$. Rest will follow from induction and a reduction to Case II.

For $k>t$, let $\lambda_k$ be the largest integers such that $(u_k' - \lambda_k u_t')\cdot \gamma >0$. Then
\begin{align*}
    \mu'_t+\sum_{k>t} \mu'_k\lambda_k &\geq \mu'_t+\sum_{k>t} \mu'_k\left(\frac{u_k'\cdot \gamma}{u_t'\cdot \gamma}-1\right) \\
    &= \frac{\sum_{k<t} (-\mu'_k)u'_k\cdot \gamma + e\cdot \gamma }{u'_t\cdot \gamma} - M \\
    &\geq \frac{e\cdot \gamma}{\delta } - M > 0 
\end{align*}

Now let 
$$z_i = \begin{cases}
    u'_i & i\leq t \\
    u'_i - \lambda_iu_t' & i>t.
\end{cases} $$

Then $Z\in \cM(U')$ and we have
\begin{align*}
    e &= \sum_k \mu_k'u'_k \\
    &= \sum_{k<t} \mu_k'u_k'+ \sum_{k>t} \mu_k'(u_k'- \lambda_k u_t') + \left(\mu'_t+\sum_{k> t}\mu_k\lambda_k\right) u_t'. \\
    & = \sum_{k\neq t} \mu_k' z_k + \theta_t z_t
\end{align*}

where $\theta_t = \sum_{k>t}\mu_k\lambda_k + \mu_t' > 0$. 
\end{proof}

\end{document}
